% Generated draft: requires mathematical and source review.
\documentclass[11pt]{article}
\usepackage[margin=25mm]{geometry}
\usepackage{fontspec}
\setmainfont{Latin Modern Roman}
\usepackage{amsmath,amssymb,mathtools,unicode-math}
\setmathfont{Latin Modern Math}
\usepackage{enumitem,needspace,xurl,hyperref}
\setlength{\parindent}{0pt}
\setlength{\parskip}{5pt}
\setlength{\emergencystretch}{3em}
\newcommand{\N}{\mathbb{N}}
\newcommand{\Z}{\mathbb{Z}}
\newcommand{\Q}{\mathbb{Q}}
\newcommand{\R}{\mathbb{R}}
\newcommand{\C}{\mathbb{C}}
\newcommand{\F}{\mathbb{F}}
\newcommand{\A}{\mathbb{A}}
\newcommand{\PP}{\mathbb{P}}
\newcommand{\E}{\mathbb{E}}
\newcommand{\eps}{\varepsilon}
\newcommand{\dd}{\,\mathrm{d}}
\begin{document}
\section*{UnsolvedMath ID: 10600026}
\textbf{Source problem code:} catalogue entry
\subsection{Definitions and mathematical statement}
Let \(F\) be a specified supporting surface and let \(K\subset F\times I\), \(I=[0,1]\), be a knot in its thickening. The usual virtual-knot interpretation uses orientable supporting surfaces; the excerpt does not state all surface or boundary conventions. In an arrow-polynomial state, a state curve \(c\subset F\) is assigned its surface isotopy class \([c]_F\) and its arrow number \(a(c)\). Retaining both gives state variables of the form \(X_{[c]_F,a(c)}\), rather than variables recording the arrow number alone. The arrow-number convention, cancellation rules, treatment of contractible curves, coefficients, and normalization are not supplied.

\emph{Investigation problem: the surface arrow invariant.} Develop and investigate this surface-dependent arrow-polynomial refinement for knots in specified thickened surfaces. At the level of notation its intended form can be expressed schematically as
\[
 P_F(K)=\sum_{s}w(s)\prod_{c\in s}X_{[c]_F,a(c)},
\]
where the sum is over the intended arrow states and \(w(s)\) denotes their unspecified weights. This schematic expression records the two required labels; it is not a fully defined state sum or an asserted new formula. A precise construction must give the state rules and prove invariance for the intended surface-knot equivalence.

Investigate its application to a virtual knot by choosing a minimal-surface representative. A minimal representative is intended to minimize supporting genus in the relevant stable surface model. To obtain a virtual-knot invariant one must specify how surface diffeomorphisms act on the isotopy-class labels and verify independence of the allowed choices of minimal representative, or prescribe the appropriate equivalence class of polynomial values. The excerpt does not provide those comparison conventions. It describes the refinement as more powerful than the ordinary arrow polynomial and very strong, and as generalizing work on the surface bracket polynomial; these qualitative assessments are preserved without a claimed complete invariant or a quantified separating theorem.
\subsection{Short English statement}
Investigate an arrow polynomial for a knot in a fixed thickened surface whose state variables remember both the surface isotopy class and arrow number of each state curve. Applying it to virtual knots via minimal surfaces also requires a precise rule for comparing labels across surface representatives.
\subsection{Sources}
\begin{itemize}
\item Local UnsolvedMath registry, permanent ID 10600026, source code catalogue entry, ``Virtual-knot problem 26 --- surface arrow invariant.'' No individual proposer is supplied.
\item \emph{Unsolved Problems in Virtual Knot Theory and Combinatorial Knot Theory}, revised survey, problem-list item 26. Supplied URL: \url{https://arxiv.org/abs/1409.2823}; recorded access 2026-07-29; extraction source-tex. The dedicated section has 60 top-level enumerated items and retains compound source items intact. The statement's citation for previous surface-bracket work is missing.
\item The supplied summary lists ``On two categorifications of the arrow polynomial for virtual knots'' (2009), with placeholder \texttt{arXiv:0906.xxxx}; ``The homological arrow polynomial for virtual links'' (2022), \texttt{arXiv:2207.xxxx}; ``On arrow polynomials of checkerboard colorable virtual links'' (2020), \texttt{arXiv:2002.xxxx}; and ``Surface pole bracket polynomials of virtual knots and twisted knots'' (2014), \texttt{arXiv:1401.xxxx}. These incomplete identifiers are retained as supplied and are not usable verified citations.
\item Supplied status partially\_solved and PARTIAL-PROGRESS. The summary describes categorifications and surface variants as existing and says the program appears partially developed, while reporting no dedicated paper found beyond those listed. This is a supplied search assessment. Current primary-literature status and redistribution terms for the publicly accessible author/publisher copy remain NEEDS\_REVIEW.
\end{itemize}
\end{document}
